% Research excerpt retained from Paper II, Sections 6.2--6.3.
% Source snapshot: 9f76211bd58bf490b31487852c908eeb7d1cb755.
% Not a standalone manuscript: uses the complex/characteristic-zero setup,
% marked normal closure, and macros of paper/paper2.tex at that snapshot.
% Retained for the broader planar-vanishing program, not the degree-two proof
% and not a characteristic-two statement. See README.md and OBJECT-BENCH.md.

\subsection{The reduced target boundary}

Inside the same marked closure, set
\[
  E=B[g(\iota(A)):g\in G]\subseteq N.
\]
Thus \(E\) is generated by all conjugate source algebras.  It is the
coordinate ring of the connected component of their fiber product
selected by these embeddings: the tensor product is regular and
\'etale over \(B\), and its selected generic factor is \(N\).
In particular, \(E/B\) is \'etale and \(\Frac(E)=N\).

The action of \(G\) on \(\Spec E\) is free.  Indeed, the scheme-theoretic
fixed locus of each \(g\ne1\) is open by unramifiedness and closed by
separatedness.  Since \(g\) acts nontrivially on \(\Frac(E)=N\), this
locus misses the generic point and is empty by connectedness.
The quotient \(\Spec E\to\Spec E^G\) is therefore a finite \'etale
\(G\)-torsor \cite{StacksFiniteFreeQuotient}.
Consequently \(\Spec E^G\to\Spec B\) is \'etale, separated, and
birational, hence an open immersion \cite{StacksBirationalEtaleOpen}.
Write its image as \(U_F\).  The complement of this dense affine open
has pure codimension one \cite{StacksAffineComplement}; since
\(B=\C[P,Q]\) is factorial, we may write
\[
  U_F=D(s),\qquad
  \mathfrak b_F:=I(\Spec B\setminus U_F)=(s),
  \qquad 0\ne s\in B\ \text{squarefree}.
\]
For an empty complement take \(s=1\).  The ideal \(\mathfrak b_F\)
depends only on \(U_F\), not on its defining equation, and
\(E^G=B[1/s]\).  It is a target-base boundary ideal, distinct from
the collision obstruction \(\Obs(F)\subseteq C_F\).

\begin{proposition}[Source-regularity criterion]
\label{prop:source-regularity-criterion}
Identify \(A\), \(B\), and \(K\) with their marked images in \(N\).
Then
\[
  A\cap K=B,\qquad E^G\cap A=B,
\]
and
\[
  E^G\subseteq A
  \quad\Longleftrightarrow\quad E^G=B
  \quad\Longleftrightarrow\quad \mathfrak b_F=B.
\]
\end{proposition}

\begin{proof}
The first identity is Wang's intersection theorem
\cite[Theorem~41(i)]{Wang1980}.  We recall the divisibility argument
used in the formalization.  An irreducible \(p\in B\) remains a
nonunit in \(A\), since the units of both polynomial rings are the
nonzero constants.  Choose an irreducible factor \(q\in A\) of \(p\).
Flatness and going-down give height at most one for \((q)\cap B\);
as this prime contains \((p)\), it equals \((p)\).  Thus divisibility
by \(p\) between elements of \(B\) is reflected from \(A\) to \(B\).
Factorization and cancellation reflect arbitrary divisibility.
Clearing a fraction in \(A\cap K\) now gives \(A\cap K=B\).
Since \(E^G\subseteq N^G=K\), the second identity follows.
It gives the first equivalence, and the open immersion
\(\Spec E^G\hookrightarrow\Spec B\) gives the second.
\end{proof}

These conditions are also equivalent to \(\Obs(F)=0\).  Indeed, if
\(\mathfrak b_F=B\), then \(E\) is finite over \(B\).  Since \(B\)
is Noetherian, the \(B\)-submodule \(\iota(A)\subseteq E\) is also finite.
Thus \(F\) is
finite \'etale; the associated covering of the simply connected space
\(\C^2\) has one sheet, so \(F\) is a polynomial automorphism and
\(\Obs(F)=0\).  Conversely, the coordinate equalities in the proof of
Proposition~\ref{prop:normal-closure-criterion} give
\(E\subseteq\iota(A)\) when \(\Obs(F)=0\), and
Proposition~\ref{prop:source-regularity-criterion} gives
\(\mathfrak b_F=B\).  This reverse implication uses the collision
equations directly, without Ax--Grothendieck.  The geometric implication
from boundary emptiness uses finite-\'etale rigidity, separately from
the Galois-rigidity argument above.

\subsection{A finite first-order criterion}

Let \(D_P,D_Q\) denote the extensions of \(\partial_P,\partial_Q\)
from \(B\) to \(K\) \cite{StacksDerivationLocalization}, and set
\[
  M_0=s^{-1}B,\qquad
  M_1=M_0+B\,D_P(1/s)+B\,D_Q(1/s)\subseteq K.
\]
The module \(M_0\) has one generator, \(M_1\) has at most three, and
\(M_1/M_0\) has at most two.  All finiteness statements here are over
\(B\), not finite-dimensionality over \(\C\).

\begin{proposition}[Finite boundary criterion]
\label{prop:finite-boundary-criterion}
For every \(0\ne s\in B\), the following conditions are equivalent:
\[
  \begin{gathered}
  M_1=M_0;\\
  D_P(M_0)\subseteq M_0,\qquad D_Q(M_0)\subseteq M_0;\\
  \partial_Ps,\partial_Qs\in(s);\\
  s\in B^\times.
  \end{gathered}
\]
For the reduced boundary equation above, these conditions say
\(U_F=\Spec B\).
\end{proposition}

\begin{proof}
Since each target derivative preserves \(B\), Leibniz reduces preservation
of \(M_0\) to its generator.  For \(j=P,Q\),
\[
  D_j(1/s)=-\frac{\partial_js}{s^2}\in s^{-1}B
  \quad\Longleftrightarrow\quad \partial_js\in(s).
\]
If \(s\) is nonconstant, every nonzero partial derivative has smaller
total degree than \(s\) and cannot be a multiple of it.  Both partials
would therefore have to vanish, forcing \(s\) to be constant in
characteristic zero, a contradiction.  Nonzero constants satisfy all
the displayed conditions.
\end{proof}

The same test has an intrinsic boundary formulation.  The conormal map
\cite{StacksConormalSequence}
\[
  \mathfrak b_F/\mathfrak b_F^2
  \longrightarrow
  \Omega_{B/\C}\otimes_B B/\mathfrak b_F
  \simeq (B/\mathfrak b_F)^2
\]
sends the generator \(s\bmod s^2\) to
\[
  (\partial_Ps\bmod s,\ \partial_Qs\bmod s).
\]
Its vanishing is equivalent to the criterion above.  The intrinsic
conormal criterion does not require a choice of the generator \(s\).
It uses the existing conormal module of the boundary, not a replacement
for \(\Obs(F)\).

The first-order test is another characterization of the same empty
boundary.  Neither field-level extension of the Keller frame nor
preservation of the localized algebra \(B[1/s]\) proves it: differentiation
may send \(1/s\) to \(-\partial_js/s^2\).
